% !PRINTED 374
\weberpage{374}
% Source-faithful Cayley configuration signs: a four-stroke star and a seven-rayed star.
\providecommand{\weberfourup}{\mathord{\tikz[baseline=-0.2ex,x=0.23em,y=0.23em,line width=0.45pt]{\draw (0,0)--(-1.05,1) (0,0)--(-0.35,1.08) (0,0)--(0.35,1.08) (0,0)--(1.05,1);}}}
\providecommand{\webersevenstar}{\mathord{\tikz[baseline=-0.35ex,x=0.21em,y=0.21em,line width=0.45pt]{\draw (0,0)--(90:1) (0,0)--(38.57:1) (0,0)--(-12.86:1) (0,0)--(-64.29:1) (0,0)--(-115.71:1) (0,0)--(-167.14:1) (0,0)--(141.43:1);}}}
%<<CONT
of these determinations is therefore 280, and we arrive at the following result:

\gesperrt{There are in all 288 Aronhold systems, whose symbols have one of the two forms}

\[
\webersevenstar;\ \Delta\,\weberfourup.
\]

From these symbols, however, one must not conclude that these systems of seven split into two different kinds. The difference between the two figures lies only in the notation. In fact, we may start from any one whatever of the complete systems of seven in order to derive the notation.

From this one easily finds the notation for the \gesperrt{Steiner} complexes, which are to be formed after one of the two following types:

\[
[1\,2][3\,4],\ [1\,3][2\,4],\ [1\,4][2\,3],\ [5\,6][7\,8],\ [5\,7][6\,8],\ [5\,8][6\,7],
\]
\[
[1\,7][1\,8],\ [2\,7][2\,8],\ [3\,7][3\,8],\ [4\,7][4\,8],\ [5\,7][5\,8],\ [6\,7][6\,8].
\]

\begin{center}
\S.\,100.

\gesperrt{Rational determination of the curve from a complete system of seven.}
\end{center}

The great significance of the \gesperrt{Aronhold} systems for the problem of the double tangents is expressed in the following two theorems:

I. \gesperrt{If for a curve of the fourth order a complete system of seven is given, then from it all the remaining double tangents can be rationally determined.}

II. \gesperrt{If seven arbitrary straight lines in a plane are given, then in general, i.\,e. when certain rational functions of the coefficients in the equations of these lines do not vanish, one can determine in a rational way the equation of a curve of the fourth order without singular point, for which the given seven lines form an Aronhold system.}

In order to prove the first theorem, it would, in view of the complete equal entitlement of the figures 1 to 7, suffice to
%CONT>>
